\renewcommand{\arraystretch}{1.5}

\section{Cronograma de ejecución}
\label{subsec:cronogramaEjecucion}

El plan de ejecución organiza las acciones previstas para alcanzar los tres objetivos específicos de Comunica, indicando los procedimientos propuestos, sus instrumentos de verificación, los protagonistas y los resultados esperados.

\vspace{1.5em}
\subsection{Plan de ejecución de Comunica -- Objetivo Específico 1}

{\footnotesize
	\begin{xltabular}{\textwidth}{|
			>{\raggedright\arraybackslash\hspace{0pt}\hsize=0.45\hsize}X|
			>{\raggedright\arraybackslash\hspace{0pt}\hsize=1.2\hsize}X|
			>{\raggedright\arraybackslash\hspace{0pt}\hsize=1.75\hsize}X|
			>{\raggedright\arraybackslash\hspace{0pt}\hsize=1.0\hsize}X|
			>{\raggedright\arraybackslash\hspace{0pt}\hsize=0.8\hsize}X|
			>{\raggedright\arraybackslash\hspace{0pt}\hsize=0.8\hsize}X|
		}
		\hline
		\textbf{Objetivos} & \textbf{Acciones} & \textbf{Descripción de las acciones} & \textbf{Instrumento de verificación} & \textbf{Protagonistas} & \textbf{Resultados esperados} \\
		\hline
		\endfirsthead
		
		\multicolumn{6}{c}{\textit{Viene de la página anterior}} \\
		\hline
		\textbf{Objetivos} & \textbf{Acciones} & \textbf{Descripción de las acciones} & \textbf{Instrumento de verificación} & \textbf{Protagonistas} & \textbf{Resultados esperados} \\
		\hline
		\endhead
		
		\hline
		\multicolumn{6}{r}{\textit{Continúa en la siguiente página...}} \\
		\endfoot
		
		\endlastfoot
		
        \hyperlink{subsec:SMS}{OE1}&Revisión documental y SMS & Revisar estudios sobre registro, clasificación y seguimiento de incidencias mediante el estudio de mapeo sistemático previsto en el nivel diagnóstico. Identificar aportes útiles para Comunica. & Protocolo y matriz del SMS; registro de estudios seleccionados. & Equipo de investigación. & Referentes y funcionalidades utilizadas. \\
		\hline
		\hyperlink{subsec:instrumento}{OE1} & Crear Encuesta Fase 1 : Diseño y estructuración & Elaborar un cuestionario dirigido a ciudadanos de Rivas para evaluar las limitaciones de los métodos tradicionales de reporte y definir qué funcionalidades son necesarias en una nueva plataforma. & Cuestionario de encuesta estructurado y validado. & Equipo de investigación. & Instrumento de recolección de datos primarios diseñado y listo para su aplicación. \\
		\hline
    	\hyperlink{subsec:instrumento2}{OE1} & Implementacion de encuesta Fase 2: Aplicación y tabulación de datos & Distribuir la encuesta en la población objetivo, recopilar las respuestas y tabular los datos para cuantificar de forma real la necesidad del sistema digital en el municipio. & Informe de resultados tabulados con gráficos estadísticos y análisis de respuestas. & Equipo de investigación y usuarios. & Datos cuantitativos que demuestran las limitaciones actuales y las preferencias de los ciudadanos con respecto a funcionalidades necesarias en un sistema. \\
		\hline
	    \hyperlink{subsec:requerimentos}{OE1}& Identificación de requerimientos & Precisar los requerimientos funcionales y no funcionales a partir de la documentación y los datos primarios obtenidos en la encuesta ciudadana. & Matriz de trazabilidad. & Equipo de investigación y desarrollo. & Requerimientos de Comunica definidos. \\
		\hline
	    \hyperlink{subsec:backlog}{OE1} & Creacion de historias de usuario y Product Backlog & Transformar los requerimientos en historias de usuario priorizadas, con criterios de aceptación y su vínculo con las funcionalidades previstas. Organización por sprints (sujeto a cambios). & Inventario de historias de usuario; Product Backlog; criterios de aceptación. & Equipo de desarrollo; responsable de planificación. & Backlog inicial coherente con los requerimientos. \\
		\hline
	\end{xltabular}
}

\newpage

\subsection*{Plan de ejecución de Comunica -- Objetivo Específico 2}

{\footnotesize
	\begin{xltabular}{\textwidth}{|
			>{\raggedright\arraybackslash\hspace{0pt}\hsize=0.45\hsize}X|
			>{\raggedright\arraybackslash\hspace{0pt}\hsize=1.2\hsize}X|
			>{\raggedright\arraybackslash\hspace{0pt}\hsize=1.75\hsize}X|
			>{\raggedright\arraybackslash\hspace{0pt}\hsize=1.0\hsize}X|
			>{\raggedright\arraybackslash\hspace{0pt}\hsize=0.8\hsize}X|
			>{\raggedright\arraybackslash\hspace{0pt}\hsize=0.8\hsize}X|
		}
		\hline
		\textbf{Objetivos} & \textbf{Acciones} & \textbf{Descripción de las acciones} & \textbf{Instrumento de verificación} & \textbf{Protagonistas} & \textbf{Resultados esperados} \\
		\hline
		\endfirsthead
		
		\multicolumn{6}{c}{\textit{Viene de la página anterior}} \\
		\hline
		\textbf{Objetivos} & \textbf{Acciones} & \textbf{Descripción de las acciones} & \textbf{Instrumento de verificación} & \textbf{Protagonistas} & \textbf{Resultados esperados} \\
		\hline
		\endhead
		
		\hline
		\multicolumn{6}{r}{\textit{Continúa en la siguiente página...}} \\
		\endfoot
		
		\endlastfoot
		
		OE2 & Arquitectura de software & Definir la separación entre frontend Angular, API REST en AdonisJS y persistencia de datos en PostgreSQL, así como los componentes de autenticación y acceso. & Documento o diagrama de arquitectura; especificación de componentes. & Equipo de desarrollo (frontend y backend). & Arquitectura técnica y responsabilidades de componentes definidas. \\
		\hline
		OE2 & Modelo de base de datos & Diseñar entidades y relaciones para usuarios, incidencias, categorías, evidencias, ubicaciones, estados e historial, considerando integridad referencial y permisos. & Diagrama entidad-relación; esquema relacional; diccionario de datos. & Responsable backend y equipo de desarrollo. & Modelo de datos preparado para respaldar el registro y seguimiento. \\
		\hline
		OE2 & Diseño de interfaces UI/UX & Diseñar la interfaz de reporte, historial, consulta de estados y panel administrativo, incluyendo la selección de ubicación mediante un mapa interactivo. & Prototipos de interfaz; flujos de navegación; revisión de usabilidad. & Responsable frontend/diseño. & Interfaces y navegación definidas para usuarios y administradores. \\
		\hline
		OE2 & Modelado del flujo de incidencias & Representar el proceso de registro, clasificación, consulta, actualización de estado y trazabilidad de las incidencias, así como la información visible según el rol. & Diagramas de flujo o casos de uso; reglas funcionales documentadas. & Equipo de análisis y desarrollo. & Flujo funcional y responsabilidades de acceso delimitados. \\
		\hline
	\end{xltabular}
}

\vspace{3em}

\subsection*{Plan de ejecución de Comunica -- Objetivo Específico 3}

{\footnotesize
	\begin{xltabular}{\textwidth}{|
			>{\raggedright\arraybackslash\hspace{0pt}\hsize=0.45\hsize}X|
			>{\raggedright\arraybackslash\hspace{0pt}\hsize=1.2\hsize}X|
			>{\raggedright\arraybackslash\hspace{0pt}\hsize=1.75\hsize}X|
			>{\raggedright\arraybackslash\hspace{0pt}\hsize=1.0\hsize}X|
			>{\raggedright\arraybackslash\hspace{0pt}\hsize=0.8\hsize}X|
			>{\raggedright\arraybackslash\hspace{0pt}\hsize=0.8\hsize}X|
		}
		\hline
		\textbf{Objetivos} & \textbf{Acciones} & \textbf{Descripción de las acciones} & \textbf{Instrumento de verificación} & \textbf{Protagonistas} & \textbf{Resultados esperados} \\
		\hline
		\endfirsthead
		
		\multicolumn{6}{c}{\textit{Viene de la página anterior}} \\
		\hline
		\textbf{Objetivos} & \textbf{Acciones} & \textbf{Descripción de las acciones} & \textbf{Instrumento de verificación} & \textbf{Protagonistas} & \textbf{Resultados esperados} \\
		\hline
		\endhead
		
		\hline
		\multicolumn{6}{r}{\textit{Continúa en la siguiente página...}} \\
		\endfoot
		
		\endlastfoot
		
		OE3 & Construcción del backend & Desarrollar la API REST con AdonisJS, validación de datos, autenticación y control de acceso basado en roles (RBAC), junto con servicios para incidencias y sus estados. & Repositorio de código; documentación de endpoints; pruebas de API. & Responsable backend. & Servicios principales disponibles para integrarse con el frontend. \\
		\hline
		OE3 & Desarrollo del frontend & Implementar en Angular formularios de reporte, consulta de historial, seguimiento de estados y pantallas administrativas; integrar la comunicación con la API. & Repositorio de frontend; capturas funcionales; casos de prueba. & Responsable frontend. & Interfaz web integrada con las funciones previstas. \\
		\hline
		OE3 & Geolocalización y evidencias & Integrar el mapa interactivo mediante Leaflet/OpenStreetMap y la asociación de coordenadas a los reportes; incorporar la carga y consulta de evidencias visuales con Cloudinary (proveedor sujeto a cambios). & Pruebas de registro geográfico y carga de archivos; registros de prueba. & Equipo frontend y backend. & Reportes con ubicación y evidencias vinculadas. \\
		\hline
		OE3 & Gestión y seguimiento de incidencias & Implementar registro, clasificación, actualización de estados, consulta del historial y permisos de administración, de acuerdo con el alcance del proyecto. & Casos de prueba funcionales; matriz de criterios de aceptación. & Equipo de desarrollo. & Flujo de atención y trazabilidad funcional en los escenarios de prueba. \\
		\hline
		OE3 & Panel administrativo y métricas & Construir indicadores y visualizaciones para consultar cantidad, estados y distribución territorial de incidencias, utilizando componentes gráficos contemplados en el proyecto. & Capturas del panel; pruebas con datos de ejemplo; verificación de indicadores. & Equipo frontend y backend. & Panel que facilite la consulta y el análisis de reportes. \\
		\hline
		OE3 & Pruebas y documentación & Ejecutar pruebas funcionales y de integración sobre los casos definidos, registrar incidencias de software y elaborar documentación técnica y guías de uso. Despliegue de demostración en Render/Netlify (proveedores y modalidad de alojamiento sujetos a cambios). & Matriz de pruebas; registro de incidencias; manuales; evidencia de despliegue si se realiza. & Equipo de investigación y desarrollo. & Versión de evaluación verificada en los escenarios disponibles y documentación inicial. \\
		\hline
	\end{xltabular}
	
   El apartado de los Resultados esperados dentro del cronogrma de ejecución describe los productos, entregables o estados concretos que se alcanzarán tras completar cada acción propuesta. Su propósito  es establecer un criterio de éxito medible para cada actividad, traduciendo las tareas de investigación, diseño y desarrollo en evidencias tangibles tales como documentos de análisis, prototipos de interfaz, bases de datos estructuradas o módulos de código funcional, En conjunto, la union de estos resultados permite evaluar el avance real del proyecto Comunica, garantizando que cada paso realizado contribuya de manera verificable al cumplimiento progresivo de sus tres objetivos específicos.
   
  \textbf{Resultados Cronograma de ejecución}  
   
\textbf{1. \hypertarget{subsec:SMS}{OE1}: Referentes tecnológicos y funcionales identificados.}

El análisis de herramientas internacionales previas (\cite{parrales2025desarrollo}; \cite{trejo2019sistema}; \cite{moreno2026diseno}, entre otros) ha permitido definir un estándar funcional base para el desarrollo de «Comunica». La evidencia documental establece que los sistemas modernos de gestión de incidencias requieren integrar las siguientes capacidades:

\begin{itemize}
	\item \textbf{Geolocalización y evidencia multimedia:} Uso de coordenadas espaciales y carga de evidencia fotográfica para otorgar precisión a los reportes ciudadanos, reduciendo los tiempos de validación en campo.
	
	\item \textbf{Gestión de estados y notificaciones:} Implementación de flujos de estado basados en buenas prácticas (ITIL) y notificaciones automáticas (números de radicado o historial), garantizando que el ciudadano conozca el avance de su reporte y reduciendo la incertidumbre.
	
	\item \textbf{Paneles de control y analítica geoespacial:} Incorporación de \textit{dashboards} centralizados que agrupen la información mediante análisis geoespacial y estadístico, permitiendo a las entidades priorizar rutas y tomar decisiones operativas basadas en mapas de concentración de incidencias.
	
	\item \textbf{Clasificación y priorización estructurada:} Segmentación de los reportes por categorías (tipos de incidencia) para organizar la recepción de datos y facilitar su canalización hacia el área competente.
	
	\item \textbf{Asignación de responsables y trazabilidad integral:} Capacidad de registrar qué funcionario asume cada caso, conservando un historial inmutable de las fechas, respuestas y acciones técnicas documentadas durante todo el ciclo de atención.
	
	\item \textbf{Mecanismos de verificación y cierre:} Inclusión de procesos de validación que aseguren la resolución real del problema antes de dar por finalizada o cerrada la incidencia en el sistema.
	
	\item \textbf{Accesibilidad:} Diseño web responsivo para facilitar el acceso desde cualquier dispositivo móvil o de escritorio y, según lo requiera el contexto institucional, la posibilidad de realizar reportes anónimos.
\end{itemize}


\textbf{2.\hypertarget{subsec:instrumento}{OE1} Instrumento de recolección de datos primarios diseñado y listo para su aplicación} 

Como resultado del plan 2 se obtuvo una encuesta la cual recopila datos demográficos y examina los obstáculos y frustraciones que enfrentan las personas al utilizar canales tradicionales de comunicación. Asimismo, valora la percepción pública sobre la falta de transparencia y seguimiento por parte de las entidades proveedoras tras emitir un aviso. Se indaga también en el grado de aceptación hacia el uso de herramientas tecnológicas y plataformas web centralizadas que faciliten la gestión y reduzcan la incertidumbre ciudadana. Finalmente, el formulario busca determinar si la creacion de sistemas digitales en tiempo real lograría fortalecer la confianza y fomentar una mejor atencion al usuario en la resolución de incidencias.

\textbf{link de encuesta:}
\textcolor{blue}{https://forms.gle/mccAsdU9PHPXU5QcA}

\textbf{3.\hypertarget{subsec:requerimentos}{OE1} Identificación de requerimientos}

A partir del análisis documental y de campo, se definieron siete requerimientos (\textbf{MR-01} al \textbf{MR-07}) que delimitan el alcance funcional y técnico de la plataforma «Comunica». En conjunto, estos requerimientos cubren cuatro dimensiones esenciales del sistema:

\begin{itemize}
	\item \textbf{Seguridad y control de acceso (MR-01):} Garantiza la autenticación de usuarios y la diferenciación de permisos según su rol operativo (ciudadanía, personal institucional y administradores).
	\item \textbf{Captura y organización territorial (MR-02 y MR-03):} Estructura la división geográfica administrativa para ubicar cada reporte y provee un mecanismo interactivo con geolocalización y evidencia multimedia.
	\item \textbf{Enrutamiento y seguimiento operativo (MR-04 y MR-05):} Automatiza la clasificación por tipología hacia la entidad competente y asegura la trazabilidad completa durante todo el ciclo de vida de la incidencia.
	\item \textbf{Analítica y calidad del sistema (MR-06 y MR-07):} Consolida indicadores de desempeño en paneles visuales, respaldados por una arquitectura que garantiza integridad transaccional, seguridad y alto rendimiento.
\end{itemize}
 Para una vista especifica de los requerimentos del proyecto revise los Metarequisitos plasmados en la \textbf{\hyperlink{subsec:matriz}{Matriz de Trazabilidad}} 
 
 \textbf{4.\hypertarget{subsec:backlog}{OE1} \hyperlink{subsec:backlog1}{Cómo resultado del plan de historias usuarios se concreto el Backlog, este se encuentra en la página 60, toque para visualizar}}
 


 

}